\section{讲座定位：开放对齐不是一个算法，而是一段系统史}

本讲不是一份 ``RLHF 101''。Nathan Lambert 要回答的是：ChatGPT 之后，开放社区如何从只会续写的 base model，走到能遵循指令、对话、拒答并接受偏好优化的模型；哪些数据、评测和训练方法真正改变了生态，哪些名称只是阶段性标签。阅读这堂课时，应该把模型、数据、算法、硬件与评测看成相互约束的系统，而不是寻找唯一“对齐配方”。

\subsection{先纠正来源，再理解课程范围}

本节先说明材料边界。Lecture 28 的官方讲者是 AI2 的 Nathan Lambert，主题是开放语言模型对齐；旧目录误用了上一讲 Hyung Won Chung 的架构历史 deck。正确来源是 77 页官方 Google Slides 和 1,693 条人工字幕，因此本讲必须整篇重写，而不是在旧文本上补丁式修改。

\lecturefigure{slide-001.jpg}{Nathan Lambert 的课程主题：Aligning Open Language Models}{official deck, p.~1}

\begin{importantbox}{本讲的核心判断}
开放对齐的进展来自一条链：\textbf{更强 base model} 提供能力，\textbf{instruction/preference data} 指定交互分布，\textbf{fine-tuning 与 RLHF/DPO} 改变输出概率，\textbf{evaluation} 提供迭代信号，\textbf{开放发布} 让社区复现并产生下一轮数据。链中任何一环不足，单个算法都无法复制 ChatGPT 体验。
\end{importantbox}

\teachervoice{Nathan 在 00:00:34--00:01:03 强调，这不是基础教程，而是把 ChatGPT 之后的 fine-tuning 与 alignment 历史重新讲一遍，让听众知道哪些概念仍重要、哪些只是短暂热潮。}

\subsection{从 Shannon 到 ChatGPT：为什么要先讲旧历史}

接下来十页是一条逐步展开的时间轴。它不是为了争论“谁最早”，而是为了说明对齐建立在语言建模、Transformer、规模化、风险批评和产品交互之上。每个里程碑都改变了后来开放社区能下载什么、能训练什么、以及需要修复什么。

\lecturefigure{slide-003.jpg}{1948：Shannon 用字符序列近似英语}{official deck, p.~3}

自回归语言模型的基础目标可写为
\[
\mathcal{L}_{\mathrm{LM}}(\theta)=-\sum_{t=1}^{T}\log p_{\theta}(x_t\mid x_{<t}),
\]
其中 $x_t$ 是第 $t$ 个 token，$x_{<t}$ 是历史 token，$\theta$ 是模型参数。该目标只要求提高真实 continuation 的概率；“有帮助”“无害”“服从系统指令”都不是它直接编码的属性。

\lecturefigure{slide-005.jpg}{2017：Transformer 让 attention 成为主干计算结构}{official deck, p.~5}

\lecturefigure{slide-006.jpg}{2018：ELMo、GPT-1 与 BERT 奠定现代预训练范式}{official deck, p.~6}

\begin{knowledgebox}{读图：三条路线后来怎样汇合}
ELMo 强调上下文化表示，BERT 强调双向编码与下游适配，GPT 强调自回归生成。今天的 chat model 主要沿生成路线发展，却吸收了表示学习、迁移学习和大规模预训练的共同经验。时间轴支持的是技术积累，而不是单一路线从一开始就必然胜出。
\end{knowledgebox}

\lecturefigure{slide-007.jpg}{2019：GPT-2、模型发布争论与 scaling law 进入主线}{official deck, p.~7}

\lecturefigure{slide-008.jpg}{GPT-2 发布策略使能力、风险与开放性同时成为问题}{official deck, p.~8}

\lecturefigure{slide-009.jpg}{2020：GPT-3 展示规模化能力，也放大偏见与伤害讨论}{official deck, p.~9}

\begin{warningbox}{能力增长与对齐不是同一条曲线}
更低 next-token loss 能提高通用生成能力，却不会自动确定回答风格、拒答边界、事实引用或用户偏好。规模化让“可被对齐的能力空间”更大，也可能让不希望的行为更强；对齐阶段是在已有分布上重新加权，而不是凭空创造全部能力。
\end{warningbox}

\lecturefigure{slide-010.jpg}{2021：Stochastic Parrots 把数据、偏见、环境与权力问题推到中心}{official deck, p.~10}

\lecturefigure{slide-011.jpg}{2022：ChatGPT 把对话式对齐变成大众可感知的产品体验}{official deck, p.~11}

\teachervoice{00:01:25--00:05:11 的时间轴把 autoregressive loss、Transformer、scaling、harms 与 ChatGPT 连成一条因果背景。讲者的重点是：不能只从 RLHF 论文开始讲对齐，因为它依赖此前二十多年的模型与基础设施。}

\subsection{ChatGPT 能否没有 RLHF}

时间轴在 ChatGPT 处停下后，课程提出一个更窄的问题：RLHF 是否必要。课堂答案是 ``necessary, but not sufficient''——偏好优化能显著改变可用性和安全行为，但产品还依赖预训练模型、数据清洗、prompt interface、system policy、工程反馈与持续评测。

\lecturefigure{slide-012.jpg}{问题：ChatGPT 能否在没有 RLHF 的情况下存在}{official deck, p.~12}

\lecturefigure{slide-014.jpg}{Anthropic 的偏好优化曲线显示 RLHF 可显著改变模型行为}{official deck, p.~14}

\lecturefigure{slide-015.jpg}{Llama 2 等公开模型也把 RLHF 作为关键训练阶段}{official deck, p.~15}

\begin{knowledgebox}{读图：跨公司的 Elo 曲线不能直接横比}
每张曲线的 prompt、标注者、比较对手和 Elo 标尺都可能不同。图能说明同一组织内部经过 RLHF 后偏好胜率改善，却不能把 Anthropic 曲线上的数值与 Llama 曲线数值当作统一排行榜。比较前必须确认评测分布与标定方式。
\end{knowledgebox}

\teachervoice{00:05:11--00:06:56，Nathan 一方面用公司内部曲线说明 RLHF 的巨大影响，另一方面马上提醒这些 Elo 没有跨组织校准。课程因此把“RLHF 重要”与“这些图能证明多少”分开处理。}

\subsection{本章小结}

开放对齐是一层建立在 base model 之上的行为塑形系统。历史页说明能力、风险与交互需求如何逐步累积；RLHF 页说明偏好训练能产生显著行为变化，却不是完整产品的充分条件。下一章将建立术语与章节地图，避免把 IFT、SFT、RLHF、DPO 和 alignment 混成同一个词。

